\chapter{Sprint 3 : Planning intelligent et assistant d'apprentissage}
\label{chap:sprint3}

\section{Introduction}
\addcontentsline{toc}{section}{Introduction}

Le troisième sprint marque un tournant dans le développement de \textit{Smart Focus} :
il s'agit du moment où le système cesse d'être une simple application de suivi pour
devenir un véritable assistant pédagogique. Les fonctionnalités développées durant cette
itération reposent toutes sur l'intelligence artificielle et visent à accompagner
l'utilisateur dans deux dimensions complémentaires de son apprentissage : l'organisation
de son temps et l'exploitation de ses ressources documentaires.

Sur le plan de la planification, le système est désormais capable de générer
automatiquement des sessions d'étude adaptées au profil de l'utilisateur, à ses examens
à venir et à son état de récupération tel que mesuré par le score de sommeil. Du côté
de l'assistance documentaire, trois outils ont été mis en place : un chatbot RAG qui
permet d'interroger directement le contenu de documents PDF, un générateur de quiz QCM
et un système de flashcards piloté par l'algorithme de répétition espacée SM-2.

Ce sprint constitue ainsi le cœur applicatif du projet, en connectant les interfaces
mobiles aux modules IA du backend et en rendant le système réellement utile au quotidien.

\section{Backlog du Sprint 3}

Le backlog du Sprint 3 regroupe les fonctionnalités liées à l'intelligence artificielle
et à l'assistance pédagogique. Ces éléments constituent le noyau applicatif du système
et ont été priorisés en fonction de leur valeur pour l'utilisateur et de leurs
interdépendances techniques — le pipeline RAG, par exemple, est le socle commun
du chatbot, des quiz et des flashcards.

Le tableau~\ref{tab:backlogs3} présente les quatre user stories du sprint avec leurs
tâches associées.

\begin{table}[H]
\centering
\renewcommand{\arraystretch}{1.3}
\begin{tabular}{|c|p{3cm}|p{6cm}|p{4cm}|}
\hline
\textbf{ID} & \textbf{Fonctionnalité} & \textbf{User Story} & \textbf{Tâche} \\
\hline

US9 & Chatbot intelligent
& En tant qu'utilisateur, je souhaite interagir avec un chatbot pour poser des
  questions et réviser mes cours
& Développement d'un assistant de révision capable de répondre aux questions
  de l'utilisateur à partir de ses propres documents \\ \hline

US10 & Import de documents
& En tant qu'utilisateur, je souhaite importer mes cours au format PDF afin
  qu'ils soient utilisés par le chatbot
& Développement de la fonctionnalité d'importation et d'indexation de documents
  pour les rendre exploitables par le système \\ \hline

US11 & Génération de quiz et flashcards
& En tant qu'utilisateur, je souhaite générer des quiz QCM et des flashcards
  pour tester et consolider mes connaissances
& Développement des outils de génération de quiz et de flashcards avec
  planification automatique des révisions \\ \hline

US12 & Planification intelligente
& En tant qu'utilisateur, je souhaite obtenir un planning de révision adapté
  à mon profil et à mon état de récupération
& Développement d'un système de planification automatique adapté au profil
  et aux disponibilités de l'utilisateur \\ \hline

\end{tabular}
\caption{Backlog du Sprint 3}
\label{tab:backlogs3}
\end{table}

\section{Spécifications fonctionnelles}

\subsection{Diagramme de cas d'utilisation du Sprint 3}

Le diagramme de cas d'utilisation ci-dessous illustre les principales interactions
entre l'utilisateur et le système durant ce sprint. Il couvre les fonctionnalités
liées à l'import de documents, au chatbot RAG, à la génération de quiz et de
flashcards, ainsi qu'à la planification intelligente.

\begin{figure}[H]
\centering
\includegraphics[width=1\textwidth]{images/Sprint3UC1.png}
\caption{Diagramme de cas d'utilisation du Sprint 3}
\end{figure}
\begin{figure}[H]
\centering
\includegraphics[width=1\textwidth]{images/Sprint3UC2.png}
\caption{Diagramme de cas d'utilisation du Sprint 3 Continue}
\end{figure}

\subsection{Description textuelle des cas d'utilisation}

Les fiches suivantes décrivent le comportement attendu du système pour chaque cas
d'utilisation du sprint. Elles précisent les conditions de déclenchement, les étapes
du scénario nominal et les situations d'erreur à gérer.

%----------------------------------------
\subsubsection{Cas d'utilisation : Uploader un document}

L'import de documents est le point d'entrée de l'ensemble des fonctionnalités d'assistance. Sans document importé, aucun des outils (chatbot, quiz, flashcards) ne peut fonctionner. Le processus va de la sélection du fichier par l'utilisateur jusqu'à sa mise à disposition complète pour les fonctionnalités du système. Le tableau~\ref{tab:uc_upload} en décrit le déroulement.

\begin{table}[H]
\centering
\renewcommand{\arraystretch}{1.2}
\begin{tabular}{|p{4cm}|p{10cm}|}
\hline
\textbf{Acteur} & Utilisateur \\ \hline
\textbf{Objectif} & Importer un document PDF pour l'exploiter via le chatbot,
                    les quiz et les flashcards \\ \hline
\textbf{Conditions initiales} & Utilisateur connecté, fichier PDF valide \\ \hline
\textbf{Déroulement} &
\begin{itemize}
    \item L'utilisateur sélectionne un fichier depuis son appareil
    \item L'utilisateur envoie le fichier au système
    \item Le système extrait le contenu textuel du document
    \item Le système découpe le contenu en segments exploitables
    \item Le système analyse et indexe chaque segment pour la recherche
    \item Le système enregistre les informations du document
\end{itemize} \\ \hline
\textbf{Cas d'erreur} & Fichier vide ou sans contenu lisible → refus avec
                        message d'erreur \\ \hline
\textbf{Résultat} & Document disponible pour le chatbot, les quiz et les flashcards \\ \hline
\end{tabular}
\caption{Description du cas d'utilisation « Uploader un document »}
\label{tab:uc_upload}
\end{table}

%----------------------------------------
\subsubsection{Cas d'utilisation : Poser une question RAG}

Une fois les documents importés, l'utilisateur peut interroger leur contenu en langage naturel. Le système identifie les passages les plus pertinents du document, construit une réponse ciblée et la présente avec les références correspondantes. Le tableau~\ref{tab:uc_rag_chat} détaille ce processus.

\begin{table}[H]
\centering
\renewcommand{\arraystretch}{1.2}
\begin{tabular}{|p{4cm}|p{10cm}|}
\hline
\textbf{Acteur} & Utilisateur \\ \hline
\textbf{Objectif} & Obtenir une réponse basée sur le contenu de ses documents importés \\ \hline
\textbf{Conditions initiales} & Au moins un document importé \\ \hline
\textbf{Déroulement} &
\begin{itemize}
    \item L'utilisateur saisit sa question dans l'interface du chatbot
    \item L'utilisateur sélectionne les documents à interroger
    \item Le système recherche les passages les plus pertinents dans les documents
    \item Le moteur d'intelligence artificielle génère une réponse à partir de ces passages
    \item L'utilisateur voit la réponse avec les références aux sources utilisées
    \item Le système enregistre l'échange
\end{itemize} \\ \hline
\textbf{Cas d'erreur} & Aucun document sélectionné → le système répond sans se
                        baser sur les documents \\ \hline
\textbf{Résultat} & Réponse avec références aux passages sources \\ \hline
\end{tabular}
\caption{Description du cas d'utilisation « Poser une question au chatbot »}
\label{tab:uc_rag_chat}
\end{table}

%----------------------------------------
\subsubsection{Cas d'utilisation : Générer un quiz}

Au-delà de la consultation passive, le système offre à l'utilisateur un moyen actif
de tester sa compréhension. À partir d'un ou plusieurs documents importés, le système
génère des questions à choix multiples accompagnées d'explications, permettant à l'utilisateur
d'évaluer ses lacunes de manière structurée. Le tableau~\ref{tab:uc_quiz} présente
ce processus.

\begin{table}[H]
\centering
\renewcommand{\arraystretch}{1.2}
\begin{tabular}{|p{4cm}|p{10cm}|}
\hline
\textbf{Acteur} & Utilisateur \\ \hline
\textbf{Objectif} & Tester ses connaissances à travers un quiz QCM généré
                    automatiquement \\ \hline
\textbf{Conditions initiales} & Au moins un document indexé ou une session
                                 terminée avec documents liés \\ \hline
\textbf{Déroulement} &
\begin{itemize}
    \item L'utilisateur sélectionne le ou les documents sources
    \item L'utilisateur choisit le nombre de questions souhaité (3 à 30)
    \item Le système analyse le contenu des documents sélectionnés
    \item Le moteur d'intelligence artificielle génère des questions à choix multiples avec explications
    \item Le système enregistre le quiz
    \item L'utilisateur voit les questions s'afficher (les bonnes réponses sont masquées avant soumission)
\end{itemize} \\ \hline
\textbf{Cas d'erreur} & Contenu insuffisant dans le document →
                        message d'erreur \\ \hline
\textbf{Résultat} & Quiz QCM prêt à être complété par l'utilisateur \\ \hline
\end{tabular}
\caption{Description du cas d'utilisation « Générer un quiz »}
\label{tab:uc_quiz}
\end{table}

%----------------------------------------
\subsubsection{Cas d'utilisation : Générer et réviser des flashcards}

Les flashcards complètent le dispositif de révision en s'appuyant sur un principe scientifiquement établi : la répétition espacée. Le système calcule dynamiquement le moment optimal pour représenter chaque carte à l'utilisateur, en fonction de la qualité de son rappel lors des sessions précédentes. Le tableau~\ref{tab:uc_flashcard} décrit ce mécanisme.

\begin{table}[H]
\centering
\renewcommand{\arraystretch}{1.2}
\begin{tabular}{|p{4cm}|p{10cm}|}
\hline
\textbf{Acteur} & Utilisateur \\ \hline
\textbf{Objectif} & Réviser ses cours à l'aide de flashcards avec répétition
                    espacée \\ \hline
\textbf{Conditions initiales} & Au moins un document importé \\ \hline
\textbf{Déroulement} &
\begin{itemize}
    \item L'utilisateur sélectionne le ou les documents sources
    \item Le moteur d'intelligence artificielle génère des cartes avec une question au recto et sa réponse au verso
    \item L'utilisateur révise une carte et évalue la facilité avec laquelle il a rappelé la réponse
    \item Le système calcule automatiquement la prochaine date de révision optimale
    \item Le système affiche uniquement les cartes dont la révision est prévue ce jour
\end{itemize} \\ \hline
\textbf{Cas d'erreur} & Document sans contenu suffisant →
                        message d'erreur \\ \hline
\textbf{Résultat} & Paquet de flashcards avec planification automatique
                    des révisions \\ \hline
\end{tabular}
\caption{Description du cas d'utilisation « Générer et réviser des flashcards »}
\label{tab:uc_flashcard}
\end{table}

%----------------------------------------
\subsubsection{Cas d'utilisation : Générer un planning intelligent}

La planification intelligente est la fonctionnalité la plus intégrée du sprint :
elle croise le profil utilisateur, les données de sommeil, les examens à venir et
l'emploi du temps importé pour produire un planning de révision personnalisé.
La génération suit une approche hybride — déterministe pour la mécanique des
créneaux, IA pour le choix des sujets — garantissant à la fois la fiabilité
et la pertinence du résultat. Le tableau~\ref{tab:uc_planning} en décrit les étapes.

\begin{table}[H]
\centering
\renewcommand{\arraystretch}{1.2}
\begin{tabular}{|p{4cm}|p{10cm}|}
\hline
\textbf{Acteur} & Utilisateur \\ \hline
\textbf{Objectif} & Obtenir un planning de révision optimisé et adapté à
                    son état \\ \hline
\textbf{Conditions initiales} & Utilisateur authentifié avec un profil
                                 configuré \\ \hline
\textbf{Déroulement} &
\begin{itemize}
    \item Le système analyse le profil de l'utilisateur et ses disponibilités
    \item Si un emploi du temps a été fourni, le système l'analyse pour identifier les créneaux libres
    \item Le système calcule les créneaux disponibles en respectant les pauses entre sessions
    \item Les sessions sont filtrées selon la préférence horaire de l'utilisateur (matin, après-midi, soir)
    \item La durée des sessions est adaptée en fonction de l'état de récupération de l'utilisateur
    \item Le moteur d'intelligence artificielle attribue les sujets à réviser à chaque créneau
    \item En cas d'indisponibilité du service IA, le système attribue les sujets automatiquement
    \item Le système enregistre les sessions dans le planning de l'utilisateur
\end{itemize} \\ \hline
\textbf{Cas d'erreur} & Aucun créneau disponible dans la journée →
                        aucune session générée \\ \hline
\textbf{Résultat} & Planning journalier avec sessions adaptées au profil
                    utilisateur \\ \hline
\end{tabular}
\caption{Description du cas d'utilisation « Générer un planning intelligent »}
\label{tab:uc_planning}
\end{table}

\section{Conception}

\subsection{Diagrammes de séquence}

Les diagrammes de séquence suivants décrivent le comportement dynamique du système
pour les principaux cas d'utilisation du sprint. Ils mettent en évidence les
échanges entre l'interface mobile, le backend FastAPI, les services IA et les
couches de persistance.

%----------------------------------------
\subsubsection{Processus d'upload et d'indexation d'un document}

Comme illustré dans la Figure~\ref{fig:seq_upload}, l'upload d'un document déclenche
une chaîne de traitements côté backend : extraction du texte par PyMuPDF \cite{pymupdf}, découpage
en chunks, vectorisation via le modèle HuggingFace \cite{sentencetransformers}, puis stockage parallèle dans
ChromaDB \cite{chromadb} pour les vecteurs et dans PostgreSQL \cite{postgresql} pour les métadonnées. C'est ce pipeline
qui rend le document exploitable par toutes les fonctionnalités IA du sprint.

\begin{figure}[H]
\centering
\includegraphics[width=0.85\textwidth]{images/uploadSeq.png}
\caption{Diagramme de séquence — upload et indexation d'un document}
\label{fig:seq_upload}
\end{figure}

%----------------------------------------
\subsubsection{Processus de chat RAG}

La Figure~\ref{fig:seq_chat_rag} illustre le parcours d'une question posée
au chatbot. La recherche sémantique dans ChromaDB \cite{chromadb} permet de sélectionner les
passages les plus pertinents du document, qui sont ensuite assemblés en contexte
avant d'être transmis à Groq \cite{groq}. La réponse retournée cite ses sources, offrant
à l'utilisateur une traçabilité complète des informations présentées.

\begin{figure}[H]
\centering
\includegraphics[width=0.85\textwidth]{images/RagSeq.png}
\caption{Diagramme de séquence — chat RAG}
\label{fig:seq_chat_rag}
\end{figure}

%----------------------------------------
\subsubsection{Processus de génération de quiz}

La Figure~\ref{fig:seq_quiz} détaille la génération d'un quiz QCM. Le backend
construit un échantillon diversifié de chunks à partir de ChromaDB, formule un
prompt structuré et soumet la requête à Groq. La réponse JSON est ensuite validée
et persistée en base de données avant d'être transmise à l'interface mobile.

\begin{figure}[H]
\centering
\includegraphics[width=0.85\textwidth]{images/quizSeq.png}
\caption{Diagramme de séquence — génération de quiz}
\label{fig:seq_quiz}
\end{figure}

%----------------------------------------
\subsubsection{Processus de génération du planning intelligent}

La Figure~\ref{fig:seq_planning} illustre l'approche hybride adoptée pour la
planification. Dans un premier temps, un algorithme déterministe calcule les
créneaux disponibles en tenant compte des blocs occupés, du buffer entre sessions
et des préférences horaires de l'utilisateur. Groq \cite{groq} intervient ensuite uniquement
pour attribuer des sujets pertinents à ces créneaux pré-calculés. En cas d'échec
du modèle, un mécanisme de fallback déterministe garantit que l'utilisateur
obtient toujours un planning valide.

\begin{figure}[H]
\centering
\includegraphics[width=0.85\textwidth]{images/PlanningSeq.png}
\caption{Diagramme de séquence — génération du planning intelligent}
\label{fig:seq_planning}
\end{figure}

\subsection{Diagramme de classes du Sprint 3}

Le diagramme de classes présenté dans la Figure~\ref{fig:class_sprint3} modélise
les entités du domaine mobilisées durant ce sprint. Il s'articule autour de trois
pôles principaux.

Le premier concerne l'assistant documentaire : \textit{ChatDocument} représente
un document PDF importé et indexé. Il est lié aux échanges du chatbot via
\textit{ChatMessage}, et constitue la source des quiz et des flashcards. La classe
\textit{Quiz} agrège une collection de \textit{QuizQuestion}, chacune portant ses
options, sa correction et son explication. La classe \textit{Flashcard} intègre
quant à elle les paramètres de l'algorithme SM-2 (\textit{ease\_factor},
\textit{interval}, \textit{repetitions}, \textit{next\_review}), qui pilotent
la planification automatique des révisions.

Le second pôle concerne la planification : \textit{StudySession} représente une
session d'étude, qu'elle soit générée par l'IA ou créée manuellement. Elle est
liée à plusieurs documents (relation many-to-many avec \textit{ChatDocument}),
peut être associée à un quiz et déclencher la génération de flashcards. La classe
\textit{Exam} permet à l'utilisateur de déclarer ses échéances, qui sont prises
en compte lors de la génération du planning.

Enfin, l'ensemble des entités est rattaché à \textit{User} et \textit{UserProfile},
qui centralisent respectivement les données de compte et les préférences personnelles
(objectif de focus quotidien, plage horaire préférée) utilisées pour personnaliser
le planning.

\begin{figure}[H]
\centering
\includegraphics[width=0.95\textwidth]{images/Sprint3ClassV2.png}
\caption{Diagramme de classes du Sprint 3}
\label{fig:class_sprint3}
\end{figure}

\subsection{Choix du modèle de langage}
\label{sec:choix_modele}

L'ensemble des fonctionnalités IA du système repose sur un modèle de langage
accessible via API. Ce choix est structurant : il conditionne la qualité des
réponses générées, la fiabilité du format JSON produit pour les quiz et le
planning, ainsi que la fluidité perçue dans l'interface de chat.

\subsubsection*{Critères de sélection}

Quatre critères ont guidé l'évaluation des alternatives disponibles :

\begin{itemize}
    \item \textbf{Vitesse d'inférence} : dans une interface de chat, la latence
    est directement perçue par l'utilisateur. Un modèle lent dégrade l'expérience
    même si ses réponses sont de bonne qualité.

    \item \textbf{Fiabilité du format JSON} : la génération de quiz et la
    planification nécessitent que le modèle retourne systématiquement une structure
    JSON valide et exploitable. Tout écart impose un fallback déterministe.

    \item \textbf{Fenêtre de contexte} : le pipeline RAG transmet plusieurs chunks
    de documents en même temps. Une fenêtre courte force à tronquer le contexte,
    dégradant la qualité des réponses.

    \item \textbf{Accessibilité} : dans le cadre d'un projet académique, la
    disponibilité d'un tier gratuit utilisable sans carte bancaire est un critère
    déterminant.
\end{itemize}

\subsubsection*{Comparaison des alternatives}

Le tableau~\ref{tab:model_comparison} confronte les principaux modèles envisagés
sur ces quatre axes, ainsi que leur adéquation au contexte du projet.

\begin{table}[H]
\centering
\renewcommand{\arraystretch}{1.4}
\begin{tabular}{|p{3.8cm}|p{2cm}|p{2cm}|p{2.2cm}|p{1.8cm}|p{2cm}|}
\hline
\textbf{Modèle} & \textbf{Vitesse} \newline (tokens/s)
                & \textbf{JSON} \newline \textbf{fiabilité}
                & \textbf{Contexte}
                & \textbf{Tier} \newline \textbf{gratuit}
                & \textbf{RPM gratuit} \\
\hline

\textbf{Groq \cite{groq} — llama-3.3-70b \cite{llama3}} \newline \textit{(retenu)}
    & $\sim$400 & Élevée & 128 K & Oui & 30 \\
\hline

Groq \cite{groq} — llama-3.1-8b \cite{llama3}
    & $\sim$900 & Moyenne & 128 K & Oui & 30 \\
\hline

Google Gemini 2.0 Flash \cite{gemini}
    & $\sim$200 & Élevée & 1 M & Oui & 15 \\
\hline

OpenAI GPT-4o-mini \cite{openai}
    & $\sim$100 & Très élevée & 128 K & Non & — \\
\hline

Anthropic Claude Haiku 3.5 \cite{anthropic}
    & $\sim$120 & Très élevée & 200 K & Non & — \\
\hline

Mistral 8x7B (Groq) \cite{mistral}
    & $\sim$350 & Moyenne & 32 K & Oui & 30 \\
\hline

\end{tabular}
\caption{Comparaison des modèles de langage évalués}
\label{tab:model_comparison}
\end{table}

\subsubsection*{Justification du choix}

Le modèle retenu est \textbf{Llama 3.3 70B Versatile}, exécuté via la plateforme
\textbf{Groq} \cite{groq}. Ce choix se justifie par la convergence de plusieurs avantages
complémentaires.

\textbf{Sur le plan de la performance}, Groq se distingue par une architecture
matérielle propriétaire (LPU — Language Processing Unit) conçue spécifiquement
pour l'inférence de modèles de langage. Là où les API concurrentes reposent sur
des GPU classiques et délivrent typiquement entre 80 et 200 tokens par seconde,
Groq atteint environ 400 tokens par seconde pour llama-3.3-70b. Concrètement,
une réponse du chatbot RAG apparaît presque instantanément, sans latence
perceptible pour l'utilisateur.

\textbf{Sur le plan de la qualité}, Llama 3.3 70B est le modèle open-source
de Meta offrant le meilleur équilibre taille/capacité au moment du développement.
Avec 70 milliards de paramètres et un entraînement optimisé pour le suivi
d'instructions, il produit des sorties JSON cohérentes — notamment à
\texttt{temperature=0.2} tel que configuré dans le système — ce qui est
indispensable pour la génération de quiz et l'assignation des sujets de planning.
Son support multilingue couvre naturellement le français académique utilisé dans
les documents des étudiants.

\textbf{Sur le plan de l'accessibilité}, Groq propose un tier gratuit de 30
requêtes par minute sans nécessiter de carte bancaire, contre 15 RPM seulement
pour Google Gemini (l'autre fournisseur intégré dans le système). Cette limite
supérieure est directement exploitée par le rate-limiter interne du projet,
configuré à 25 RPM pour rester sous le seuil avec une marge de sécurité.

\textbf{Limites identifiées} : GPT-4o-mini et Claude Haiku 3.5 offriraient une
fiabilité JSON légèrement supérieure pour des schémas très complexes, mais à un
coût unitaire incompatible avec le cadre académique du projet. Gemini 2.0 Flash
constitue une alternative viable — raison pour laquelle le système l'intègre
comme second fournisseur configurable — mais sa limite de 15 RPM le rend moins
adapté aux phases de test intensif. Llama 3.1 8B, bien que plus rapide encore,
produit des sorties JSON insuffisamment fiables pour les cas d'usage critiques
du projet.

En définitive, \textbf{llama-3.3-70b-versatile sur Groq représente le meilleur
compromis vitesse / qualité / gratuité} pour les besoins spécifiques du système
\textit{Smart Focus}.

\section{Environnement technique}

Cette section présente les principaux outils et technologies mis en œuvre dans le cadre du Sprint 3. Ces choix ont été guidés par des critères de performance, de disponibilité et d'adéquation aux besoins pédagogiques du système.

%----------------------------------------
\subsection{Pipeline RAG (Retrieval-Augmented Generation)}

Le pipeline RAG est l'architecture centrale qui sous-tend le chatbot, les quiz et les flashcards. Il repose sur le principe suivant : plutôt que de s'appuyer uniquement sur la connaissance générale du modèle de langage, le système récupère d'abord les passages les plus pertinents du document de l'utilisateur, puis les transmet au modèle comme contexte pour générer une réponse ancrée dans le contenu réel du cours.

Ce pipeline se décompose en quatre étapes techniques : l'extraction du texte brut à partir du fichier, le découpage en segments, la transformation de ces segments en représentations numériques (vecteurs), et enfin la recherche par similarité pour identifier les passages pertinents à une question donnée.

%----------------------------------------
\subsection{PyMuPDF — Extraction du texte}

PyMuPDF \cite{pymupdf} est une bibliothèque Python permettant d'extraire le contenu textuel de fichiers PDF page par page, avec un support des documents scannés et des mises en page complexes. Elle constitue la première étape du pipeline : sans extraction fiable du texte, aucune indexation n'est possible.

Dans le système, PyMuPDF parcourt chaque page du document importé, extrait le texte brut, et transmet ce contenu à l'étape suivante.

%----------------------------------------
\subsection{LangChain — Découpage en segments}

LangChain \cite{langchain} est un framework Python facilitant la construction de pipelines fondés sur des modèles de langage. Dans ce sprint, seul le composant \texttt{RecursiveCharacterTextSplitter} est utilisé. Il découpe le texte extrait en segments de taille contrôlée (500 caractères par défaut avec un chevauchement de 50 caractères), garantissant que chaque segment soit à la fois suffisamment court pour être traité efficacement et suffisamment large pour conserver son sens.

%----------------------------------------
\subsection{Sentence Transformers (HuggingFace) — Représentation vectorielle}

Le modèle \texttt{all-MiniLM-L6-v2} \cite{sentencetransformers}, fourni par HuggingFace, transforme chaque segment textuel en un vecteur numérique de 384 dimensions qui encode son sens sémantique. Deux segments traitant du même sujet produiront des vecteurs proches, même s'ils utilisent des mots différents. C'est ce mécanisme qui permet au système de retrouver les passages pertinents à une question, sans se limiter à une correspondance mot-à-mot.

Ce modèle fonctionne entièrement en local, sans appel à un service externe, ce qui garantit la confidentialité des documents importés par l'utilisateur.

%----------------------------------------
\subsection{ChromaDB — Base de données vectorielle}

ChromaDB \cite{chromadb} est une base de données spécialisée dans le stockage et la recherche de vecteurs. Elle reçoit les segments vectorisés et les organise par collection (une collection par document importé). Lors d'une question, ChromaDB calcule la similarité cosinus entre le vecteur de la question et ceux des segments stockés, et retourne les \textit{k} passages les plus proches.

ChromaDB fonctionne en mode local (fichiers sur disque), sans nécessiter de serveur supplémentaire.

%----------------------------------------
\subsection{Algorithme SM-2 — Répétition espacée}

L'algorithme SM-2 \cite{sm2}, développé par Piotr Woźniak, est le fondement des systèmes de révision par répétition espacée. Son principe : l'intervalle avant la prochaine révision d'une carte est calculé dynamiquement en fonction de la note de rappel donnée par l'utilisateur (de 0 à 5).

Trois paramètres sont mis à jour après chaque révision :
\begin{itemize}
    \item \textbf{Facteur de facilité} (\textit{ease factor}) : ajusté selon la note de rappel, il conditionne la progression de l'intervalle.
    \item \textbf{Intervalle} : durée en jours avant la prochaine révision. Il augmente exponentiellement pour les cartes bien maîtrisées.
    \item \textbf{Répétitions} : remis à zéro si la note est inférieure à 3 (carte non maîtrisée), garantissant un repassage rapide.
\end{itemize}

Une carte avec une note de rappel élevée sera présentée à nouveau dans plusieurs semaines ; une carte mal rappelée sera représentée le lendemain.

%----------------------------------------
\subsection{Groq API — Modèle de langage}

La plateforme Groq \cite{groq} fournit un accès par API au modèle \textbf{Llama 3.3 70B Versatile} \cite{llama3}. Ce modèle est utilisé pour toutes les tâches de génération du sprint : réponses du chatbot, questions de quiz et sujets du planning. Le détail des critères de sélection et la comparaison avec les alternatives disponibles sont présentés dans la section~\ref{sec:choix_modele}.

%----------------------------------------
\section{Implémentation}

\subsection{Interfaces de l'application mobile}

Les interfaces développées dans le cadre du Sprint 3 donnent accès aux modules
d'intelligence artificielle du système. Elles ont été conçues pour rester
simples d'utilisation tout en exposant des fonctionnalités sophistiquées :
interrogation de documents, évaluation de connaissances et planification adaptative.

%----------------------------------------
\subsubsection{Interface du chatbot RAG}

L'interface du chatbot permet à l'utilisateur de sélectionner ses documents sources
et d'engager une conversation contextualisée. Les réponses générées indiquent les
pages de référence utilisées, renforçant la confiance de l'utilisateur dans les
informations fournies. L'historique des échanges est conservé et consultable
à tout moment.

\begin{figure}[H]
\centering
\includegraphics[width=0.35\textwidth]{images/chatbot1.jpg}
\hspace{0.5cm}
\includegraphics[width=0.35\textwidth]{images/chatbot2.jpg}
\caption{Interface du chatbot RAG : gestion des documents et conversation}
\label{fig:chatbot_screens}
\end{figure}

%----------------------------------------
\subsubsection{Interfaces de génération et passage de quiz}

L'interface de quiz guide l'utilisateur depuis la sélection des documents jusqu'aux
résultats détaillés. Chaque question est présentée avec ses quatre options et,
après soumission, l'explication de la bonne réponse est affichée pour favoriser
la compréhension plutôt que la simple mémorisation du résultat.

\begin{figure}[H]
\centering
\includegraphics[width=0.35\textwidth]{images/quiz1.jpg}
\hspace{0.5cm}
\includegraphics[width=0.35\textwidth]{images/quiz2.jpg}
\caption{Interfaces de quiz : génération et passage du quiz}
\label{fig:quiz_screens}
\end{figure}

%----------------------------------------
\subsubsection{Interfaces des flashcards avec révision espacée}

L'interface de flashcards adopte un format recto/verso classique. Après avoir
retourné une carte, l'utilisateur évalue son niveau de rappel sur une échelle
de 0 à 5. Cette note est immédiatement traitée par l'algorithme SM-2, qui
recalcule l'intervalle avant la prochaine révision de cette carte.

\begin{figure}[H]
\centering
\includegraphics[width=0.35\textwidth]{images/flashcard1.jpg}
\hspace{0.5cm}
\includegraphics[width=0.35\textwidth]{images/flashcard2.jpg}
\caption{Interfaces des flashcards : génération et révision}
\label{fig:flashcard_screens}
\end{figure}

%----------------------------------------
\subsubsection{Interface du planning intelligent}

L'interface de planning présente les sessions générées sous forme de vue
quotidienne ou hebdomadaire. L'utilisateur peut visualiser d'un coup d'œil
les blocs de révision attribués par le système, consulter le sujet et la
priorité de chaque session, et compléter le planning par des entrées manuelles
si nécessaire.

\begin{figure}[H]
\centering
\includegraphics[width=0.35\textwidth]{images/planning1.jpg}
\hspace{0.5cm}
\includegraphics[width=0.35\textwidth]{images/planning2.jpg}
\caption{Interface du planning intelligent : vue quotidienne et génération IA}
\label{fig:planning_screens}
\end{figure}

\section{Conclusion}
\addcontentsline{toc}{section}{Conclusion}

À l'issue de ce troisième sprint, \textit{Smart Focus} dispose d'un ensemble
cohérent de fonctionnalités pédagogiques pilotées par l'intelligence artificielle.
Le chatbot RAG, les quiz QCM, les flashcards SM-2 et la planification adaptative
forment un écosystème dans lequel chaque outil se nourrit des mêmes documents
importés par l'utilisateur, garantissant une expérience unifiée.

L'architecture retenue pour la planification — calcul déterministe des créneaux
combiné à une personnalisation par Groq \cite{groq} avec fallback automatique — illustre
une approche pragmatique : l'IA enrichit le résultat sans en compromettre la
fiabilité. Ce principe guidera également les développements du sprint suivant,
consacré au monitoring comportemental en temps réel et aux alertes intelligentes.